% Time form: Present tense

\chapter{Methodology}\label{ch:methodology}

This chapter defines the evaluation strategy used to assess the agentic \gls{AI} integration across the three research questions.
The evaluation is structured around a controlled benchmark in a reproducible virtual environment and a real-world assessment in the BFH network security lab.
A static reference scenario serves as the baseline against which all \gls{AI}-driven configurations are compared.

\section{Evaluation Setup}\label{sec:evaluation-setup}

Three models are benchmarked across 10 runs in three environments (\emph{containerlab}, \emph{physical lab}, and \emph{BFH Cyber Lab}):
a frontier model (\mintinline{bash}{Claude Opus 4.7}),
an institutionally hosted open-weight model (\mintinline{bash}{gpt-oss:120b}),
and a locally-run open-source model (\mintinline{bash}{qwen3:30b} via Ollama).

\subsection{Criteria}\label{subsec:criteria}

The evaluation criteria are divided into quantitative and qualitative categories.
Quantitative criteria are collected automatically by the benchmarking suite introduced in Section~\ref{sec:benchmark-suite}.
Qualitative criteria require manual assessment and are scored against the scale defined in Table~\ref{tab:qualitative-scale}.

\textbf{Qualitative criteria:}

\begin{itemize}
    \item \textbf{Correctness and Task completion/completeness:} Are the scenarios completed fully and are the results correct?
    \item \textbf{Consistency network discovery results:} How much does the output vary across $n = 10$ runs with an identical prompt, per model and environment?
\end{itemize}

\textbf{Quantitative criteria:}

\begin{itemize}
    \item \textbf{Speed:} Time from scenario start to complete output of all discovered hosts and services
    \item \textbf{Token usage: AI-specific} Number of tokens consumed per run (\gls{AI} Reconnaissance only)
    \item \textbf{Success factor:} Number of valid runs with structured output and findings.
    The factor is presented in the form of (e.g.\ 10/12 means 10 usable runs out of 12 attempts).

\end{itemize}

To fix most of the influencing factors we defined a system prompt and a scenario prompt to set the starting line for all models equally.
With 3 models $\times$ 3 environments and 1 implemented static scenario $\times$ $n = 10$ runs, this yields 120 measurements for Reconnaissance.

\subsection{Environments}\label{subsec:environments}

Three environments are used across the evaluation, each representing a different level of control, reproducibility and realism.

\begin{enumerate}
    \item Virtual environment (containerlab)
    \item Physical environment
    \item BFH lab
\end{enumerate}

\subsection{Models}\label{subsec:models}

Table~\ref{tab:used-inference-models} lists the \glspl{LLM} used and how they were accessed.
The three models span a range of parameter counts, deployment modes, and capability levels,
allowing the evaluation to draw conclusions about the suitability of different model categories for agentic network reconnaissance.

\begin{table}[h]
    \centering
    \begin{tabularx}{\linewidth}{lX}
        \hline
        \textbf{Model} & \textbf{Provider} \\
        \hline
        qwen3:30b & Self-hosted via \gls{Ollama} \\
        gpt-oss:120b & BFH-TI Machine Learning Management Platform \\
        claude-opus-4-7 & \gls{Anthropic} \gls{API} \\
        \hline
    \end{tabularx}
    \caption{Evaluated Inference Models}
    \label{tab:used-inference-models}
\end{table}

\begin{table}[h]
    \centering
    \caption{Qualitative Evaluation Scale for LLM Reconnaissance Results}
    \label{tab:qualitative-scale}
    \begin{tabularx}{\textwidth}{@{}lllX@{}}
        \toprule
        \textbf{Criterion} & \textbf{Score} & \textbf{Level} & \textbf{Description} \\
        \midrule
        \multirow{5}{*}{\textbf{Correctness}}
        & 1--2  & Insufficient  & Results are mostly wrong or irrelevant; no actionable information \\
        & 3--4  & Poor          & Few correct findings; many errors in services, ports, or versions \\
        & 5--6  & Adequate      & Majority of hosts/services correct; isolated errors in version or protocol details \\
        & 7--8  & Good          & Findings largely correct and complete; minor inaccuracies \\
        & 9--10 & Excellent     & All key hosts, services, and versions correctly and precisely identified \\
        \midrule
        \multirow{5}{*}{\textbf{Hallucination}}
        & 9--10  & Severe        & Numerous fabricated hosts, services, or vulnerabilities without basis \\
        & 7--8  & Noticeable    & Several unsubstantiated findings; confusion of IPs or services \\
        & 5--6  & Moderate      & Isolated unconfirmed statements; version numbers or banners slightly distorted \\
        & 3--4  & Low           & Rare inaccuracies; findings are traceable to raw scan data \\
        & 1--2 & None          & All statements directly verifiable from scan output; no fabrications \\
        \bottomrule
    \end{tabularx}
\end{table}
